
\chapter{Basic of EM Waves}

\section{Maxwell's Equations}

\section{In Time Domain}

Electromagnetic phenomena described by \textbf{Maxwell equations}.

\begin{align*}
\nabla \times \bar{e} &= - \frac{\partial \bar{b}}{\partial t} & \nabla \cdot \bar{d} &= \rho \\[1em]
\nabla \times \bar{h} &= \bar{J} + \frac{\partial \bar{d}}{\partial t} & \nabla \cdot \bar{b} &= 0
\end{align*}

where:

\begin{center}
\renewcommand{\arraystretch}{1.5}
\begin{tabular}{l l l l}
$\bar{e}(\bar{r}, t)$ & is & \textbf{ELECTRIC FIELD} & $[V/m]$ \\
$\bar{d}(\bar{r}, t)$ & is & the electric flux density & $[C/m^2]$ \\
$\bar{h}(\bar{r}, t)$ & is & \textbf{MAGNETIC FIELD} & $[A/m]$ \\
$\bar{b}(\bar{r}, t)$ & is & \textbf{MAGNETIC FLUX DENSITY} & $[T = Wb/m^2]$ \\
$\rho(\bar{r}, t)$ & is & \textbf{CHARGE DENSITY} & $[C/m^3]$ \\
$\bar{J}(\bar{r}, t)$ & is & \textbf{ELECTRIC CURRENT DENSITY} & $[A/m^2]$ \\
$\bar{r} = (x, y, z)$ & & & \\
\end{tabular}
\end{center}

\section{Algebra Remainder}

\begin{equation*}
\nabla \quad \text{is \textbf{NABLA}} \implies \nabla = 
\begin{bmatrix}
\frac{d}{dx} \\[0.5em]
\frac{d}{dy} \\[0.5em]
\frac{d}{dz}
\end{bmatrix}
\end{equation*}

Linear product:
\begin{align*}
\bar{a} \cdot \bar{b} &= a_x b_x + a_y c_y + a_z c_z \\
\nabla \cdot d &= \frac{\partial d_x}{\partial x} + \frac{\partial d_y}{\partial y} + \frac{\partial d_z}{\partial z}
\end{align*}


% PAGINA 06
% Uso un box con bordo rosso ma testo nero per la sezione "Important"
\begin{tcolorbox}[enhanced, colback=white, colframe=red, title=\textbf{Important:}, coltitle=red, sharp corners]
\textbf{$\nabla \cdot (\nabla \times \bar{A}) = 0$, the curl of a vector field is by definition Solenoidal (no divergence)}

Same applies when $\bar{a} \cdot (\bar{a} \times \bar{b}) = 0$ (\textit{\small ORTHOGONAL!})
\end{tcolorbox}

\section{Harmonic Domain}

\begin{itemize}
    \item In many cases, we can limit the study of EM fields to sinusoidal (or harmonic) signals.
    \item Convenient to represent fields by their phasor using the so called \textbf{complex Steinmetz vectors}.
\end{itemize}

\begin{equation*}
\bar{e}(\bar{r}, t) = \text{Re}\left[ \bar{E}(\bar{r}) \exp(j\omega t) \right]
\end{equation*}

In this course: \textbf{NO ELECTRIC CHARGES} or \textbf{ELECTRIC CURRENTS}.

\begin{align*}
\nabla \times \bar{E} &= -j \omega \bar{B} & \nabla \cdot \bar{D} &= 0 \\[1em]
\nabla \times \bar{H} &= j \omega \bar{D} & \nabla \cdot \bar{B} &= 0
\end{align*}

\section{Constitutive Relations}

The relations between fields and flux densities depend on the medium (mezzo) where the magnetic field is present.

3 media types: \textbf{LINEAR}, \textbf{ISOTROPIC}, \textbf{PIECEWISE HOMOGENEOUS}.

The constitutive relations are:
\begin{equation*}
\bar{D} = \varepsilon \bar{E} \quad \quad \bar{B} = \mu \bar{H}
\end{equation*}

where:
\begin{itemize}
    \item $\varepsilon(\omega)$ is the electric permeability $[F/m]$ (Farad/m)
    \item $\mu(\omega)$ is the magnetic permeability $[H/m]$ (Henry/m)
\end{itemize}


% PAGINA 07
In vacuum: parameters related to speed of light $\implies$ constant by definition.

\begin{align*}
c_0 &= 299\,792\,458 \ m/s & \mu_0 &= 4\pi \times 10^{-7} \ H/m
\end{align*}

\begin{equation*}
c_0 = \frac{1}{\sqrt{\mu_0 \varepsilon_0}} \implies \varepsilon_0 = \frac{1}{\mu_0 c_0^2} \approx 8,85 \cdot 10^{-12} \ F/m
\end{equation*}

For generic medium: reflective index defined as $n = \sqrt{\frac{\mu \varepsilon}{\mu_0 \varepsilon_0}}$.

\section{Helmholtz Equations}

In a linear, isotropic piecewise homogeneous medium without electric charges and currents, Maxwell's equations read:

\begin{align*}
\nabla \times \bar{E} &= -j\omega \mu \bar{H} & \nabla \cdot \bar{E} &= 0 \\[1em]
\nabla \times \bar{H} &= j\omega \varepsilon \bar{E} & \nabla \cdot \bar{H} &= 0
\end{align*}

% Derivation steps
\begin{align*}
1) \quad & \nabla \times (\nabla \times \bar{E}) = \nabla \times (-j\omega \mu \bar{H}) \\
2) \quad & \nabla \times (\nabla \times \bar{E}) = -j\omega \mu (\nabla \times \bar{H}) \\
3) \quad & \nabla \times (\nabla \times \bar{E}) = -j\omega \mu (j\omega \varepsilon \bar{E}) \\
4) \quad & \nabla \times (\nabla \times \bar{E}) = \omega^2 \mu \varepsilon \bar{E} \\
5) \quad & \nabla \times \nabla \times \bar{E} = \nabla(\nabla \cdot \bar{E}) - \nabla^2 E \quad (\textit{Note: } \nabla \cdot \bar{E} = 0) \\
6) \quad & \nabla \times \nabla \times \bar{E} = \nabla(0) - \nabla^2 E = - \nabla^2 E \\
7) \quad & -\nabla^2 E = \omega^2 \mu \varepsilon E \implies \nabla^2 E + \omega^2 \mu \varepsilon E = 0
\end{align*}

% Risultato finale evidenziato in un box giallo come nell'originale
\begin{yellowbox}
\begin{equation*}
\nabla^2 \bar{E} + k^2 \bar{E} = 0 \quad \quad \nabla^2 \bar{H} + k^2 \bar{H} = 0
\end{equation*}
\end{yellowbox}

where
\begin{equation*}
k = \omega \sqrt{\varepsilon \mu}
\end{equation*}